\chapter{Chagas Disease (American Trypanosomiasis)}

\section*{Definition and Overview}
Chagas disease, also known as American trypanosomiasis, is a parasitic infection caused by the protozoan \textit{Trypanosoma cruzi}. It is primarily transmitted to humans through the faeces of blood-sucking insects known as triatomine bugs, or "kissing bugs", which typically live in cracks and crevices of poorly constructed rural dwellings. The illness occurs in two distinct phases. The acute phase follows soon after infection and is frequently mild or completely asymptomatic. The chronic phase may develop years or even decades later and most commonly involves the heart, causing an inflammatory cardiomyopathy, and less often the digestive tract, producing oesophageal and colonic dilatation. Chagas disease is a major public health problem across Latin America, but because infected people can migrate and the infection can be passed on through blood products, organ transplantation, and pregnancy, it is now also encountered worldwide.

\section*{Epidemiology}
An estimated 6--7 million people worldwide are infected with \textit{T. cruzi}, and a further 70 million are considered at risk of acquiring the infection. The disease is endemic across 21 countries in Latin America, where rural vector-borne transmission was historically dominant. Thanks to successful vector-control campaigns, housing improvements, and systematic blood-donor screening, the number of new infections has fallen markedly over recent decades, and some countries have been certified as having interrupted transmission. Nevertheless, the burden of chronic disease among people infected years ago remains substantial. Because of international migration, autochthonous and imported cases are increasingly reported in the United States, Canada, Europe, Japan, and Australia. In non-endemic regions, transmission occurs chiefly through congenital spread, blood transfusion, and organ transplantation rather than through the insect vector.

\section*{Aetiology and Pathophysiology}
\medskip
\textbf{Transmission routes.} The parasite is transmitted in several ways:
\begin{itemize}
    \item \textbf{Vector-borne (most common in endemic areas):} the triatomine bug bites the skin and defecates nearby; the parasite is released in the faeces and enters the body through the bite wound, breaks in the skin, or the conjunctiva
    \item \textbf{Congenital:} an infected mother can transmit the parasite to her baby during pregnancy
    \item \textbf{Blood and organ transfer:} transfusion of infected blood, or transplantation of organs from infected donors
    \item \textbf{Oral transmission:} contamination of food or drink with crushed bugs or their faeces, a recognised cause of outbreaks
    \item \textbf{Accidental exposure:} rare laboratory or occupational exposure
\end{itemize}

\medskip
\textbf{Pathophysiology.} In the acute phase, \textit{T. cruzi} invades local cells and disseminates through the bloodstream, multiplying inside cells of many tissues. The immune response usually controls the parasitaemia, and symptoms resolve, but the infection is not eliminated. The parasite persists within cells of the heart, smooth muscle, and autonomic ganglia, where low-level inflammation and immune-mediated damage continue over many years. In the chronic phase this persistent inflammatory injury produces fibrosis and scarring of the myocardium, damage to the conduction system, and enlargement of the cardiac chambers. Similar damage to the oesophagus and colon can disrupt coordinated muscle contraction, leading to dilatation known as megaoesophagus and megacolon.

\section*{Clinical Features}
\medskip
\textbf{Acute phase.} Often mild or subclinical. When symptoms occur, they begin roughly one to two weeks after exposure and include:
\begin{itemize}
    \item Fever, fatigue, headache, and generalised body aches
    \item \textbf{Roma\~na sign:} painless swelling of one eyelid, typical when the parasite enters through the conjunctiva
    \item \textbf{Chagoma:} a swollen, red, indurated nodule at the site of the bite
    \item Enlargement of the liver and spleen, and occasionally lymph node swelling
    \item In infants, fever, poor feeding, and irritability
\end{itemize}

\medskip
\textbf{Chronic phase.} Most infected people remain asymptomatic for years or decades. When complications develop, they are predominantly cardiac or gastrointestinal:
\begin{itemize}
    \item \textbf{Cardiac:} palpitations, chest discomfort, breathlessness on exertion, ankle swelling, dizziness, and syncope
    \item \textbf{Digestive:} difficulty swallowing, regurgitation, heartburn, chronic constipation, and abdominal bloating
    \item \textbf{Advanced disease:} features of heart failure, stroke, or sudden cardiac death due to malignant arrhythmias
\end{itemize}

\section*{Differential Diagnosis}
\begin{itemize}
    \item Other acute febrile illnesses in tropical and travel settings (malaria, dengue, typhoid, leptospirosis, and acute HIV)
    \item Viral myocarditis, ischaemic heart disease, and other causes of dilated cardiomyopathy
    \item Idiopathic or familial dilated cardiomyopathy
    \item Achalasia, oesophageal cancer, and other causes of dysphagia and oesophageal dilatation
    \item Chronic constipation and megacolon due to other causes
    \item Periorbital oedema from allergy, cellulitis, or insect bites in the acute phase
    \item Valvular heart disease and heart failure from other causes in the chronic phase
\end{itemize}

\section*{When to Seek Medical Care}
Anyone who has lived in or travelled to an endemic area and develops fever, eyelid swelling, palpitations, chest pain, or digestive symptoms should be assessed by a doctor. Testing is particularly important for babies born to infected mothers, and for recipients of blood products or organs from possibly infected donors. Even people without symptoms should be tested if they have had a recognised exposure.

\medskip
\textbf{Call 000 (or local emergency services) for:}
\begin{itemize}
    \item Chest pain, fainting, or a very fast or irregular heartbeat
    \item Severe breathlessness or new swelling of the legs or abdomen
    \item Sudden weakness, facial droop, or difficulty speaking, suggesting stroke
    \item An infant who is unwell, feverish, or not feeding normally
\end{itemize}

\section*{Diagnosis and Investigations}
\begin{itemize}
    \item \textbf{History:} travel or residence in endemic areas, housing conditions, blood transfusions, organ transplants, and family history of the infection
    \item \textbf{Acute phase testing:} microscopic examination of fresh blood for motile parasites, or polymerase chain reaction (PCR) testing
    \item \textbf{Chronic phase testing:} serological tests for antibodies to \textit{T. cruzi}; because no single test is fully reliable, diagnosis requires at least two different positive tests
    \item \textbf{Electrocardiogram (ECG):} to detect conduction abnormalities and arrhythmias, including right bundle branch block and ventricular arrhythmias
    \item \textbf{Echocardiogram:} to assess chamber size, wall motion, and pumping function
    \item \textbf{Ambulatory rhythm monitoring:} for unexplained palpitations or syncope
    \item \textbf{Barium swallow and oesophageal manometry:} for swallowing difficulties
    \item \textbf{Barium enema and colon studies:} for severe constipation or suspected megacolon
    \item \textbf{Screening of contacts:} siblings, children of infected mothers, and household members
\end{itemize}

\section*{Management}
\begin{itemize}
    \item \textbf{Antiparasitic therapy:} benznidazole or nifurtimox are used to eliminate the parasite; they are most effective in the acute phase, in children, and in recently acquired or congenital infection
    \item \textbf{Chronic infection:} antiparasitic treatment may be offered to adults without advanced cardiomyopathy, as it can slow disease progression; decisions are individualised after specialist assessment
    \item \textbf{Congenital infection:} treatment of affected newborns is highly effective and should begin early
    \item \textbf{Heart failure:} standard therapy including ACE inhibitors or angiotensin receptor blockers, beta-blockers, and diuretics
    \item \textbf{Arrhythmia management:} anti-arrhythmic drugs, and in selected cases an implantable cardioverter-defibrillator or pacemaker
    \item \textbf{Anticoagulation:} for people at high risk of stroke, particularly those with an apical aneurysm or significant left ventricular dysfunction
    \item \textbf{Digestive complications:} dietary measures, drugs to improve swallowing and bowel transit, and dilatation or surgery for severe megaoesophagus or megacolon
    \item \textbf{Follow-up:} lifelong monitoring of cardiac and digestive function, as disease can progress silently
\end{itemize}

\section*{Complications}
\begin{itemize}
    \item Chronic Chagas cardiomyopathy -- progressive dilatation, fibrosis, and impaired pumping of the heart
    \item Heart failure and its sequelae
    \item Malignant ventricular arrhythmias and sudden cardiac death
    \item Stroke and systemic embolism from thrombus forming in the damaged heart
    \item Megaoesophagus, causing dysphagia, regurgitation, and aspiration pneumonia
    \item Megacolon, causing severe constipation, obstruction, and volvulus
    \item Malnutrition and weight loss from advanced digestive disease
    \item Congenital transmission from an infected mother to her baby
    \item Adverse effects of antiparasitic drugs, including rash, peripheral neuropathy, and gastrointestinal upset
\end{itemize}

\section*{Prognosis and Outcomes}
Many infected people remain entirely well for life, and it is estimated that roughly a third of chronically infected individuals eventually develop cardiac, digestive, or combined complications, often decades after the initial infection. The most serious outcome is Chagas cardiomyopathy, which can progress to heart failure, arrhythmias, and sudden death. Early antiparasitic treatment in the acute phase, in children, and in congenital infection can cure the infection or substantially reduce complications. For those with established chronic disease, regular cardiology and gastroenterology surveillance permits early detection and management of complications, improving long-term outcomes. Prognosis is individual and should be discussed with specialists in infectious diseases, cardiology, and gastroenterology.

\section*{Prevention and Self-Care}
\begin{itemize}
    \item Repair and improve housing in endemic areas to eliminate cracks in walls and roofs where bugs nest
    \item Use residual insecticide spraying and insecticide-treated bed nets in infested dwellings
    \item Use insect repellent and sleep under bed nets in high-risk rural areas
    \item Screen blood and organ donors so that transfusion and transplant transmission is prevented
    \item Test pregnant women from endemic areas and their newborns so congenital infection is treated promptly
    \item Wash and cook food carefully and avoid potentially contaminated beverages in endemic regions
    \item People with known infection should attend regular cardiac and digestive reviews
    \item Avoid donating blood if you have had Chagas disease or a possible exposure
\end{itemize}

\section*{Sources and Further Reading}
The following reputable sources provide further information on Chagas disease:
\begin{itemize}
    \item World Health Organization. \textit{Chagas disease fact sheets and technical guidance.} https://www.who.int/
    \item Centers for Disease Control and Prevention. \textit{Chagas disease resources for health professionals.} https://www.cdc.gov/
    \item NHS. \textit{Health A--Z: information on tropical and infectious diseases.} https://www.nhs.uk/conditions/
    \item Mayo Clinic. \textit{Diseases and conditions library.} https://www.mayoclinic.org/
    \item MSD Manual. \textit{Professional edition: parasitic infections.} https://www.msdmanuals.com/
    \item MedlinePlus. \textit{Health topics on parasitic infections.} https://medlineplus.gov/
\end{itemize}
